% Original PDF: source/普通高考/2014/2014北京文.pdf, page 2.
% PDF embedded-image attempt: pdfimages -list found no image objects; source is vector artwork.
% Grid evidence: tmp/redraw_2014_beijing_liberal/page_grid100.jpg; q20_orig.png
% was cropped from the no-grid page render after coarse-grid location.
% Elements: frequency-density y-axis and x-axis arrows, nine contiguous width-2 bars,
% x ticks 0--18, labels 频率/组距 and 阅读时间, and dashed guides at a and b.
% Relations: group intervals [0,2),...,[16,18) have frequencies
% 6, 8, 17, 22, 25, 12, 6, 2, 2 from the source table; dividing by 100 and
% group width 2 gives bar heights 0.03, 0.04, 0.085, 0.11, 0.125, 0.06,
% 0.03, 0.01, 0.01. Thus a=0.085 (top of [4,6)) and b=0.125 (top of [8,10)).
% solid segments: bar outlines, axes, tick marks, and arrowheads.
% dashed segments: guides from the y-axis to x=4 at a and x=10 at b.
% labels: a and b align with their guides; axis labels and ticks stay clear of bars.
\begin{tikzpicture}[baseline]
\begin{axis}[
    width=8.0cm,
    height=6.0cm,
    axis lines=none,
    clip=false,
    xmin=-1.25,
    xmax=20.8,
    ymin=-0.018,
    ymax=0.145,
    ticks=none,
]
    % Dashed levels a and b end at the corresponding bar boundaries.
    \draw[dashed,line width=0.55pt] (axis cs:-0.02,0.085) -- (axis cs:4,0.085);
    \draw[dashed,line width=0.55pt] (axis cs:-0.02,0.125) -- (axis cs:10,0.125);

    % Frequency densities computed from the source frequency table.
    \addplot[
        ybar interval,
        draw=black,
        fill=white,
        line width=0.65pt,
    ] coordinates {
        (0,0.03)
        (2,0.04)
        (4,0.085)
        (6,0.11)
        (8,0.125)
        (10,0.06)
        (12,0.03)
        (14,0.01)
        (16,0.01)
        (18,0)
    };

    \draw[->,line width=0.75pt] (axis cs:0,0) -- (axis cs:20.5,0);
    \draw[->,line width=0.75pt] (axis cs:0,0) -- (axis cs:0,0.14);
    \node[anchor=south west,inner sep=0pt,font=\small,align=center] at (axis cs:0.45,0.137) {\(\dfrac{\text{频率}}{\text{组距}}\)};
    \node[anchor=west,inner sep=0pt,font=\small] at (axis cs:16.0,0.018) {阅读时间};
    \frequencyylabel[-4pt]{0}{0.085}{\(a\)}
    \frequencyylabel[-4pt]{0}{0.125}{\(b\)}

    \draw[line width=0.45pt] (axis cs:0,0) -- (axis cs:0,-0.007);
    \draw[line width=0.45pt] (axis cs:2,0) -- (axis cs:2,-0.007);
    \draw[line width=0.45pt] (axis cs:4,0) -- (axis cs:4,-0.007);
    \draw[line width=0.45pt] (axis cs:6,0) -- (axis cs:6,-0.007);
    \draw[line width=0.45pt] (axis cs:8,0) -- (axis cs:8,-0.007);
    \draw[line width=0.45pt] (axis cs:10,0) -- (axis cs:10,-0.007);
    \draw[line width=0.45pt] (axis cs:12,0) -- (axis cs:12,-0.007);
    \draw[line width=0.45pt] (axis cs:14,0) -- (axis cs:14,-0.007);
    \draw[line width=0.45pt] (axis cs:16,0) -- (axis cs:16,-0.007);
    \draw[line width=0.45pt] (axis cs:18,0) -- (axis cs:18,-0.007);
    \node[anchor=north,inner sep=0pt,font=\small] at (axis cs:0,-0.0075) {0};
    \node[anchor=north,inner sep=0pt,font=\small] at (axis cs:2,-0.0075) {2};
    \node[anchor=north,inner sep=0pt,font=\small] at (axis cs:4,-0.0075) {4};
    \node[anchor=north,inner sep=0pt,font=\small] at (axis cs:6,-0.0075) {6};
    \node[anchor=north,inner sep=0pt,font=\small] at (axis cs:8,-0.0075) {8};
    \node[anchor=north,inner sep=0pt,font=\small] at (axis cs:10,-0.0075) {10};
    \node[anchor=north,inner sep=0pt,font=\small] at (axis cs:12,-0.0075) {12};
    \node[anchor=north,inner sep=0pt,font=\small] at (axis cs:14,-0.0075) {14};
    \node[anchor=north,inner sep=0pt,font=\small] at (axis cs:16,-0.0075) {16};
    \node[anchor=north,inner sep=0pt,font=\small] at (axis cs:18,-0.0075) {18};
\end{axis}
\end{tikzpicture}
